
\subsubsection{3.1 Design Rationale}

The experiment controls the primary confound in prior studies: information volume. Rather than comparing static feedback versus execution feedback (which varies content), the design compares static-to-execution versus execution-to-static (which varies only order).

Both conditions receive identical feedback content; only presentation order differs.

\subsubsection{3.2 Experimental Conditions}

\textbf{Condition A (Static-to-Execution):}
\begin{itemize}
\item Static analysis feedback (500 tokens)
\item Execution feedback (500 tokens)
\end{itemize}

\textbf{Condition B (Execution-to-Static):}
\begin{itemize}
\item Execution feedback (500 tokens)
\item Static analysis feedback (500 tokens)
\end{itemize}

Feedback content is generated once per problem-iteration pair, then presented in both orders. Deterministic truncation (first 500 tokens per type) ensures identical content.

\subsubsection{3.3 Feedback Generation}

\textbf{Static Analysis:} Pylint and Mypy are run on generated code. Output is concatenated and truncated to 500 tokens.

\textbf{Execution Feedback:} Code is executed against test cases in a sandboxed environment. Stdout/stderr, pass/fail status, and stack traces for failures are captured and truncated to 500 tokens.

\subsubsection{3.4 Repair Loop}

Both conditions use identical generation/execution logic; only concatenation order differs. The loop terminates early if all tests pass or after 3 iterations.

\subsubsection{3.5 Evaluation Metrics}

\textbf{Primary Metric:} pass@1, the fraction of problems where final code passes all tests.

\textbf{Relative Improvement:} (pass@1\_A - pass@1\_B) / pass@1\_B $\times$ 100\%

\textbf{Statistical Tests:} Bootstrap CI (10,000 resamples) for 95\% CI; McNemar's test for paired comparison.

\textbf{Mechanism Metrics:}
\begin{itemize}
\item $\Delta Pass_{1}\rightarrow _2$: change in pass rate from iteration 1 to 2
\item Regression $Rate_{1}\rightarrow _2$: P(pass@iter1 ∧ fail@iter2)
\end{itemize}

